\section{Refinement by retaining complete coefficient states}
\label{sec:refinement-state}

The public operation \wl{SeriesRefine[s,h]} reuses an ordinary inverse
when its stored forward model remains sufficient at the requested cutoff.
The implementation in \path{Kernel/RefinementState.wl} supports the
individual Lagrange and Newton methods. Its state is stored in the returned
object; there is no global cache, and refinement leaves the source object
unchanged. Coordinate wrappers, grouped coefficient generation, and inputs
whose forward expansion needs additional terms use the existing source
replay path.

\paragraph{Complete Lagrange prefixes.}
Write the normalized inverse observable as
\[
 z^r\sum_{\omega<H}z^\omega Q_\omega(\log z),
 \qquad \omega=\boldsymbol{k}\cdot\boldsymbol{d},\quad
 \boldsymbol{k}\in\mathbb N^m,\quad d_j>0.
\]
The retained state contains a downward closed set of processed indices,
its one-step outside boundary, and the complete coefficient polynomial at
every processed weight. Zero polynomials are omitted from the displayed
blocks but their indices remain processed. Advancing the state takes
\emph{all} boundary indices of the smallest weight simultaneously. Every
predecessor of such an index has smaller weight, because the gaps are
positive. Consequently, no as yet undiscovered index can have the same
weight, and resonant contributions are combined before the new block is
used or displayed.

An inverse produced with a term goal may already retain this state. For
an ordinary inverse constructed at an explicit cutoff, the complete
stored blocks instead seed the state. Enumerating
$\{\boldsymbol{k}:\boldsymbol{k}\cdot\boldsymbol{d}<H\}$ and its boundary
recovers the index bookkeeping without reevaluating the previously
computed Euler coefficients. The aggregate blocks remain authoritative;
the seed does not attempt to split a resonant coefficient polynomial into
its individual contributions.

The first omitted block is also evaluated by complete weight layers.
These frontier coefficients are retained for later refinement, including
zero layers skipped after cancellation. Caching a frontier can require
more indices than merely evaluating that frontier for a remainder. If
the additional state would exceed \wl{"MaxTerms"}, the implementation
performs the ordinary independent frontier calculation without caching
it. Thus this optimization does not reject an otherwise admissible
requested cutoff solely because it tried to retain an extra frontier.
As in ordinary construction, the bounded search through cancelled
frontier layers may return a conservative remainder at the first possible
omitted weight.

\paragraph{A bounded cache of logarithmic polynomial powers.}
The Lagrange state also retains powers of each nonconstant logarithmic
coefficient polynomial. Exponents $0$ and $1$ refer to $1$ and the original
model coefficient; larger cached powers form an initial consecutive
segment. A new power is obtained by multiplying the last retained power
by the original polynomial. Thus the product entering an Euler coefficient
can reuse polynomial powers requested by earlier indices or earlier
refinement calls. The Euler differential operators and normalization
remain unchanged, and the ordinary uncached coefficient routine provides
an independent exact comparison.

Cached powers use only unused \wl{"MaxTerms"} slots after the processed
index region and its boundary. When the region grows, the largest retained
powers are evicted until their entry count fits the remaining capacity.
If no slot is available, the requested coefficient product is evaluated
directly. Eviction never discards an inverse coefficient or changes the
remainder calculation. In particular, optional caching cannot turn an
otherwise admissible index region into a resource failure. This budget
counts retained algebraic objects; it is not a byte bound on a polynomial
whose individual coefficients may themselves be large exact expressions.

\paragraph{Recovering and extending a Newton state.}
Suppose the displayed normalized observable is
$(1+U)^r$ with $r\ne0$. Its constant coefficient is $1$, so the formal
power with exponent $1/r$ uniquely recovers $1+U$ in the filtered unit
algebra. This recovers an initial Newton state from an existing observable,
including inverse powers arising from a source endpoint at infinity.
The recovered unit is checked by substituting it into the normalized
implicit equation: the residual must vanish exactly below the claimed
precision. Only then is the state admitted for further Newton steps.

If a Newton state is already present, refinement retains its unit jet,
precision, and recorded step cutoffs. Each new step increases the
precision to the smaller of twice the old precision and the requested
precision. The existing exact residual check applies at every step.
After a coarser display is requested, a stronger retained state remains
available to subsequent calls. The formal branch is selected by constant
term $1$ throughout; no numerical root selection enters this recovery.

\paragraph{Model identity and input precision.}
Retained states are bound to the normalized model, observable power,
assumptions, original function, and source endpoint and direction. A
stored signature and direct comparisons detect incompatible states, and
the displayed prefix must agree with the corresponding state prefix.
This consistency check prevents accidental reuse after changing a model;
it is not a certificate for arbitrarily edited internal associations.

A declared forward error remains a hard cap on usable inverse precision.
For a leading exponent $p$ and internal observable exponent $r$, a forward
error $\OO(u^\rho |\log u|^k)$ permits target cutoff at most
\[
 \frac{\rho-p+r}{|p|}.
\]
The remainder's logarithmic degree is combined with that transported
error at equal power. An error introduced by automatic expansion of an
analytic source is distinguished from a declared error: if refinement
requires more source terms, the constructor is replayed and the old
model state is not reused. No cached inverse coefficients justify going
beyond an explicit input-error declaration.

\paragraph{What the counters establish.}
Each refined object exposes \wl{"RefinementStatistics"} and appends it to
\wl{"RefinementHistory"}. The Lagrange counters have distinct meanings:
\begin{description}[style=nextline,leftmargin=0pt,labelwidth=0pt,labelsep=0pt,font=\normalfont]
\item[\wl{"ReusedCoefficientEvaluations"}]
Available individual index contributions. For a seed reconstructed from
stored complete blocks this is the size of the old processed region,
not work performed during seeding.
\item[\wl{"NewCoefficientEvaluations"}]
Actual new Euler coefficient evaluations, including any uncached frontier
work. Supplying a cached polynomial product still counts as one coefficient
evaluation.
\item[\wl{"ReusedBlocks"}]
Old nonzero blocks retained below the requested cutoff.
The full old display has its own count,
\wl{"AvailableSourceBlocks"}.
\item[\wl{"ReusedPolynomialPowerRequests"}]
Requests satisfied by already retained polynomial powers during this call.
\item[\wl{"NewPolynomialPowerEvaluations"}]
New expanded polynomial-power computations, including uncached requests
when the optional cache is full. This count is separate from coefficient
evaluations and excludes constant coefficients. It covers incremental
coefficient generation; an independent uncached frontier calculation is
counted separately at the coefficient level.
\item[\wl{"RetainedPolynomialPowers"}]
The number of powers of degree at least $2$ retained in the final state.
Capacity and eviction counters record how the resource budget was applied.
\end{description}
Newton statistics separately report the old and new precision, the new
step cutoffs, and Euler work used for the remainder frontier. These are
exact bookkeeping quantities; they are not a count of all arithmetic
operations and do not prove a runtime bound.

The reproducible driver \path{Tests/BenchmarkRefinement.wl} compares a
sequence of fresh public inverse constructions with one construction
followed by public refinements. It covers dense rational gaps, two and
three irrational gaps, resonant weights, and high logarithmic degrees,
using both Lagrange and Newton examples. A deep one-gap case reaches target
cutoff $64$ and also compares with the independent Catalan-number formula
for the inverse of $x+x^2$. Acceptance requires equality of
every complete block and remainder at every requested cutoff, together
with the retained-state invariants. The driver reports elapsed times and
state counters. It measures retained expression and state sizes with
\wl{ByteCount} and evaluation peaks with \wl{MaxMemoryUsed[expr]}; these
kernel measurements are not total-process memory bounds. No
machine-dependent speed or memory threshold determines acceptance.
